\chapter{Pattern and Implicit Rules}\label{ch:mk-patterns}

\autoref{sec:mk-first} ended with three complaints. This chapter answers two of
them: how to write one rule that covers every source file, and how to stop
maintaining header dependencies by hand.

\section{The Built-in Rule Database}\label{sec:mk-implicit}

Make ships knowing how to build some things already.

\begin{shell}
$ ls
main.c
$ make main.o                    # no Makefile at all
cc    -c -o main.o main.c
\end{shell}

\begin{definition}[Implicit rule]
	A rule make already knows, matching targets by their suffix. Roughly a hundred
	are built in, covering C, C++, Fortran, Pascal, \TeX{}, Yacc, Lex, and archive
	members.
\end{definition}

\begin{keys}[0.30]
	\cmd{make -p}            & Print the entire database, rules and variables   \\
	\cmd{make -p -f /dev/null} & Print only the built-ins, with no Makefile     \\
	\cmd{make -r}            & Disable them all                                 \\
\end{keys}

The C rule is worth reading, because every Makefile in this chapter is a
variation on it:

\begin{makecode}
%.o: %.c
    $(CC) $(CPPFLAGS) $(CFLAGS) -c -o $@ $<
\end{makecode}

\begin{note}
	That is why \autoref{sec:mk-conv} insisted on the conventional variable names.
	Setting \texttt{CC} and \texttt{CFLAGS} alone changes how every \texttt{.o} is
	built, and you never write the rule. It is also why \cmd{-lm} belongs in
	\texttt{LDLIBS}: the built-in link rule puts \texttt{LDLIBS} last.
\end{note}

\begin{remark}
	\cmd{make -r} (or \cmd{MAKEFLAGS += -r} inside the file) turns the database
	off. Large projects do this: searching a hundred irrelevant rules for every
	target costs measurable time, and an accidental match --- make cheerfully
	building \file{foo} from a \file{foo.c} you forgot about --- is worse than no
	match at all.
\end{remark}

\section{Pattern Rules}\label{sec:mk-pattern}

\begin{definition}[Pattern rule]
	A rule whose target contains a \texttt{\%}. The \texttt{\%} matches any
	non-empty string --- the \vocab{stem} --- and the same stem is substituted into
	the prerequisites.
\end{definition}

\begin{makecode}
%.o: %.c
    $(CC) $(CPPFLAGS) $(CFLAGS) -c -o $@ $<
\end{makecode}

Asked for \file{parser.o}, make matches the stem \texttt{parser}, looks for
\file{parser.c}, and runs the recipe with \texttt{\$@} as \file{parser.o},
\texttt{\$<} as \file{parser.c} and \texttt{\$*} as \texttt{parser}.

\begin{eg}
	One pattern rule replaces all three of the \texttt{.o} rules in
	\autoref{sec:mk-first}, and the whole Makefile becomes:
\begin{makecode}
CC       := gcc
CPPFLAGS := -Iinclude
CFLAGS   := -Wall -Wextra -std=c17 -g
SRC      := $(wildcard *.c)
OBJ      := $(SRC:.c=.o)

all: app

app: $(OBJ)
    $(CC) $(LDFLAGS) -o $@ $^ $(LDLIBS)

%.o: %.c
    $(CC) $(CPPFLAGS) $(CFLAGS) -c -o $@ $<

clean:
    rm -f app $(OBJ)

.PHONY: all clean
\end{makecode}
	Fifteen lines, and it handles any number of source files.
\end{eg}

Patterns are not restricted to compilation:

\begin{makecode}
%.pdf: %.tex
    latexmk -xelatex -interaction=nonstopmode $<

%.png: %.svg
    inkscape --export-type=png --export-filename=$@ $<

%.csv: %.json
    jq -r '.[] | [.name, .value] | @csv' $< > $@
\end{makecode}

\begin{note}
	A pattern rule is used only when make needs the target and the prerequisite
	exists (or can itself be built). Several matching patterns are resolved by
	choosing the one with the shortest stem. If nothing matches, make reports
	\texttt{No rule to make target} --- which nearly always means a typo in a
	filename rather than a missing rule.
\end{note}

\begin{moral}
	Write pattern rules in terms of the automatic variables of
	\autoref{sec:mk-auto}, never in terms of concrete filenames. \texttt{\$@} and
	\texttt{\$<} are the only way a single rule can serve every file.
\end{moral}

\section{Static Pattern Rules}\label{sec:mk-static}

A plain pattern rule applies to \emph{anything} matching it. A static pattern
rule applies only to a list you name.

\begin{center}
	\ttfamily\large
	\{targets\}: \{target-pattern\}: \{prereq-patterns\}
\end{center}

\begin{makecode}
OBJ := main.o parser.o util.o

$(OBJ): %.o: %.c
    $(CC) $(CPPFLAGS) $(CFLAGS) -c -o $@ $<
\end{makecode}

\begin{note}
	The difference matters when two kinds of file share a suffix. With a plain
	\texttt{\%.o: \%.c}, make will try to build \emph{any} \texttt{.o} it hears
	about that way, including ones that should come from assembly or from a
	generator. Restricting the rule to a known list turns a silent wrong build into
	an honest \texttt{No rule to make target}.
\end{note}

\begin{eg}
	Static patterns are also how you build one list two ways:
\begin{makecode}
RELEASE_OBJ := $(SRC:%.c=build/release/%.o)
DEBUG_OBJ   := $(SRC:%.c=build/debug/%.o)

$(RELEASE_OBJ): build/release/%.o: %.c
    $(CC) -O2 -DNDEBUG -c -o $@ $<

$(DEBUG_OBJ): build/debug/%.o: %.c
    $(CC) -O0 -g -c -o $@ $<
\end{makecode}
	Two sets of objects from one set of sources, with no risk of the flags getting
	mixed --- which, as \autoref{sec:mk-override} warned, is exactly what
	target-specific variables cannot guarantee.
\end{eg}

\section{\texorpdfstring{\texttt{VPATH} and Out-of-Tree Builds}{VPATH and Out-of-Tree Builds}}\label{sec:mk-vpath}

Real projects do not keep sources, objects and the binary in one directory.

\begin{keys}[0.28]
	\texttt{VPATH = src:include} & A colon-separated list of directories to search for \emph{any} prerequisite \\
	\texttt{vpath \%.c src}      & The same, but only for files matching a pattern    \\
	\texttt{vpath \%.h include}  & \dots{} and headers somewhere else                 \\
\end{keys}

\begin{makecode}
vpath %.c src
vpath %.h include

BUILD := build
SRC   := $(wildcard src/*.c)
OBJ   := $(patsubst src/%.c,$(BUILD)/%.o,$(SRC))

$(BUILD)/%.o: %.c
    $(CC) $(CPPFLAGS) $(CFLAGS) -c -o $@ $<
\end{makecode}

\begin{note}
	\cmd{vpath} with a pattern is almost always preferable to \texttt{VPATH}.
	\texttt{VPATH} searches those directories for everything, so a stale
	\file{build/parser.o} sitting in a searched directory can satisfy a
	prerequisite you meant to rebuild. Being specific about which patterns come
	from where avoids a whole class of confusing builds.
\end{note}

\begin{note}[The directory must exist]
	\cmd{gcc -o build/main.o} fails if \file{build/} does not exist, and make will
	not create it for you. The clean fix is an order-only prerequisite, which
	\autoref{sec:mk-orderonly} covers.
\end{note}

\section{Automatic Dependency Generation}\label{sec:mk-depgen}

This section fixes the last complaint from \autoref{sec:mk-first}, and it is the
most valuable thing in the chapter.

The problem: \file{main.c} does \cmd{\#include "parser.h"}, so editing that
header must rebuild \file{main.o}. Make cannot know this --- it does not read
your source. Written by hand, the dependency is correct on the day you write it
and wrong within a month.

The solution: the compiler already parses the includes, so ask it.

\begin{keys}[0.20]
	\cmd{-MMD} & Write a dependency file alongside each object, listing the headers it included. \cmd{-MD} also lists system headers \\
	\cmd{-MP}  & Add a dummy target for each header, so that \emph{deleting} a header does not break the build \\
\end{keys}

\begin{makecode}
CPPFLAGS += -MMD -MP
OBJ      := $(SRC:.c=.o)
DEP      := $(OBJ:.o=.d)

app: $(OBJ)
    $(CC) $(LDFLAGS) -o $@ $^ $(LDLIBS)

clean:
    rm -f app $(OBJ) $(DEP)

-include $(DEP)          # the leading dash: do not complain on the first build
\end{makecode}

Compiling \file{main.c} now also writes \file{main.d}:

\begin{makecode}
main.o: main.c parser.h util.h
parser.h:
util.h:
\end{makecode}

which is a Makefile fragment, and \cmd{-include} pulls it back in next time.
The dependency graph now maintains itself.

\begin{note}
	The leading dash in \cmd{-include} means ``do not fail if these do not exist''.
	On a clean tree there are no \texttt{.d} files yet, and without the dash the
	very first build stops with an error. The empty \texttt{parser.h:} rules come
	from \cmd{-MP} and exist so that a header you delete is treated as a target
	with no recipe rather than a missing prerequisite.
\end{note}

\begin{moral}
	Three lines --- \cmd{-MMD -MP}, a \texttt{DEP} list, and \cmd{-include} --- and
	you never write a header dependency again. Any C or C++ Makefile without them
	is one refactor away from a build that is quietly wrong.
\end{moral}

\begin{tryit}
	Add the three lines to the Makefile from \autoref{sec:mk-pattern}, build, then
	\cmd{cat main.d}. Now add an \cmd{\#include} to \file{main.c}, rebuild, and
	look again --- the new header is there, and touching it rebuilds
	\file{main.o}. Nothing in the Makefile mentions it.
\end{tryit}
